% !TeX program = pdflatex
%
% memory-paper.tex
%
% Build:  pdflatex memory-paper.tex   (run 3x: ToC, then cross-refs)
%

\documentclass[11pt,a4paper]{article}

\usepackage[T1]{fontenc}
\usepackage[utf8]{inputenc}
\usepackage[margin=2.4cm]{geometry}
\usepackage{amsmath}
\usepackage{amssymb}
\usepackage{booktabs}
\usepackage{longtable}
\usepackage{array}
\usepackage{xcolor}
\usepackage{enumitem}
\usepackage{fancyhdr}
\usepackage[hidelinks]{hyperref}

\definecolor{warn}{RGB}{160,60,0}
\definecolor{ok}{RGB}{20,110,60}
\newcommand{\bad}[1]{\textcolor{warn}{#1}}
\newcommand{\good}[1]{\textcolor{ok}{#1}}
\newcommand{\code}[1]{\texttt{\small #1}}

\pagestyle{fancy}
\fancyhf{}
\lhead{\small xT Memory Architecture}
\rhead{\small\thepage}
\renewcommand{\headrulewidth}{0.4pt}

\title{\textbf{Local First, Private by Design:}\\[0.3em]
\large The Memory Architecture of xT}
\author{Starla Cortes}
\date{August 2026}

\begin{document}
\maketitle

\begin{abstract}
\noindent
We present the memory architecture of xT, a mobile application for hormone
replacement therapy monitoring---the full persistence stack, state-management
framework, and privacy model that together determine where user data lives,
how it is protected, and what guarantees attach to it.

The architecture rests on a six-layer persistence system: SwiftData with
CloudKit for structured longitudinal data; encrypted UserDefaults for
health-identifiable configuration; the Apple Keychain for the AES-256-GCM
encryption key; in-memory caches with documented invalidation strategies for
derived state; file-system resources for baked geometry data; and
UNUserNotificationCenter for system-level reminder state. Six singleton
managers, all \code{@MainActor @Observable} and injected via SwiftUI
Environment, each own their domain exclusively and persist via a
\code{didSet}-driven pattern that eliminates save-to-crash windows.

We describe the cold-launch sequence from container creation through
manager initialization and derived-state computation; document every
persistence key with its encryption status and rationale; walk through the
encryption model from key generation in the Keychain through AES-GCM
envelopes on UserDefaults values to the automatic migration of legacy
unencrypted data; and validate the architecture against eight persistence
scenarios that exercise cold launch, migration, sync conflicts, storage
failure, developer-safety mechanisms, and encryption-key loss.

The paper is written to be checkable: every persistence decision is cited to
its source in the codebase, every encryption boundary is named, and every
limitation---from the one-way data export to the key-loss-on-restore
problem---is stated plainly.

\textbf{Keywords:} mobile health, data persistence, encryption, privacy,
SwiftData, state management, iOS security, CloudKit
\end{abstract}

\tableofcontents
\clearpage

% ======================================================================
\section{Introduction}
\label{sec:intro}

A mobile health application that records medication history, blood test
results, injection-site locations, journal entries, and user profile data
must answer a set of questions that are simple to state and difficult to
satisfy simultaneously: Where does each piece of data live? How is it
protected? What happens when the app restarts, when the device is lost, when
the user upgrades their phone, when iCloud sync conflicts, when the
encryption key is evicted from the Keychain?

xT answers these questions with a persistence architecture that is the
subject of this paper. The application's user experience and pharmacokinetic
model are described in companion papers~\cite{xtux,xtpk}; this paper concerns
the \emph{memory} that holds them together---the six-layer persistence stack,
the six-manager state framework, the encryption model, and the privacy
architecture that together determine what the application remembers, what it
protects, and what it deliberately forgets.

\subsection{The Persistence Problem}
\label{sec:persistence_problem}

Health-tracking applications face constraints that general-purpose apps do
not. A dose record entered in January must survive an app update in March,
an OS upgrade in June, and a device migration in December. A medication type
and ester combination directly identifies a user's medical regimen and must
be encrypted at rest. A blood test calibration factor derived from the user's
own lab results must be applied to every subsequent prediction without the
user needing to understand it. None of this data can leave the user's
control---the app has no server, no analytics, no third-party SDKs---which
means the persistence architecture must be self-contained within the device
and the user's own iCloud account.

These constraints produce a set of architectural requirements that shape
every persistence decision in the application:

\begin{enumerate}[nosep,leftmargin=1.4em]
\item \textbf{Durability.} Data must survive app restarts, OS upgrades, and
  device migrations without user intervention.
\item \textbf{Encryption at rest.} Health-identifiable data must be encrypted
  such that even a forensic image of the device file system does not expose
  the user's medication regimen.
\item \textbf{No server dependency.} All data lives on-device or in the
  user's iCloud account. No xT server exists, so no server can be
  compromised.
\item \textbf{Cross-device sync.} Users with multiple devices (iPhone, iPad)
  should see consistent data without manual export/import.
\item \textbf{Developer safety.} Developer tooling (sample profiles, debug
  data seeding) must never destroy real user data.
\item \textbf{Honest degradation.} When storage fails---corrupt SwiftData
  store, evicted Keychain key---the app must degrade gracefully rather than
  crashing or silently corrupting data.
\end{enumerate}

\subsection{The Architecture Response}
\label{sec:architecture_response}

xT's memory architecture addresses these requirements through three
structural decisions that operate at every layer:

\textbf{Tiered storage by data sensitivity.} Data is placed at the least
permissive tier that can serve its access pattern. Structured longitudinal
data (dose records, journal entries, blood tests) lives in SwiftData with
CloudKit sync. Health-identifiable configuration (medication type, ester,
dose schedule) lives in UserDefaults encrypted with AES-256-GCM. The
encryption key lives in the Keychain with
\code{kSecAttrAccessibleWhenUnlockedThisDeviceOnly}. Non-sensitive
preferences (reminder times, UI toggles) live in plain UserDefaults. No
single store holds all data; no store holds data it should not.

\textbf{Singleton managers as single sources of truth.} Six managers, each an
\code{@MainActor @Observable} singleton, each owning its domain exclusively.
Mutation triggers persistence through a \code{didSet} pattern, not a periodic
save loop. Views observe manager state through the SwiftUI Environment; they
never touch storage directly. The manager is the API; the store is an
implementation detail.

\textbf{Defense in depth, not defense in breadth.} The encryption key is
generated on-device, stored in the Keychain, and never transmitted. The
encrypted values in UserDefaults are opaque blobs without the key. The
Keychain accessibility attribute prevents access when the device is locked
and prevents the key from being included in unencrypted backups. Each layer
protects against a different threat; no single-layer failure compromises all
data.

\subsection{Contributions}
\label{sec:contributions}

\begin{enumerate}[nosep,leftmargin=1.4em]
\item A complete walkthrough of the six-layer persistence stack, with
  code-level citations for every store and retrieval path, from the
  \code{ModelContainer} creation in \code{xTApp.swift} to the
  \code{DataExporter}'s multi-section CSV.
\item A documented encryption model: AES-256-GCM key generation in the
  Keychain, selective encryption of UserDefaults values, and automatic
  migration of legacy unencrypted data---with the explicit boundary between
  what is encrypted and what is not, and why.
\item A manager-architecture pattern (six \code{@MainActor @Observable}
  singletons, Environment-injected) with documented rehydration,
  invalidation, and migration paths, including the \code{ProfileSnapshot}
  symbolic-swap mechanism for developer safety.
\item An in-memory caching strategy document: which derived values are
  cached (\code{heatMapCache}, \code{screeningLog}), which are deliberately
  not memoized (\code{calibrationFactor}), and how the
  \code{suppressRefresh}/\code{performBatch} guard prevents quadratic
  recomputation during bulk operations.
\item An honest accounting of memory-architecture limitations: the
  symmetric-key-loss problem, the one-way data export, the split-persistence
  trade-off, the absence of integrity verification, and the structural gaps
  that would require architectural changes to close.
\end{enumerate}

% ======================================================================
\section{Design Principles}
\label{sec:principles}

This section describes the five principles that govern every persistence
decision in xT. They are the memory analog of the UX design principles
described in the companion paper~\cite{xtux}: every store, every key, every
encryption boundary is answerable to at least one of them.

\subsection{Defense in Depth}
\label{sec:defense_in_depth}

Every piece of data in xT is stored at the least-permissive tier that can
serve its access pattern. The six tiers are:

\begin{table}[h]
\centering
\small
\begin{tabular}{@{}p{2.2cm}p{3.2cm}p{2.4cm}p{6.8cm}@{}}
\toprule
\textbf{Tier} & \textbf{Technology} & \textbf{Accessibility} & \textbf{Data Types} \\
\midrule
1. Encryption Key & Keychain & When unlocked, this device only &
   AES-256-GCM symmetric key \\
\midrule
2. Encrypted Config & UserDefaults + AES-GCM & When unlocked &
   Medication type, ester, dose dates, profile name, profile photo \\
\midrule
3. Structured Data & SwiftData + CloudKit & Any time (CloudKit encrypted) &
   Dose records, journal entries, blood tests, screenings, preferences \\
\midrule
4. Plain Config & UserDefaults & Any time &
   Reminder times, UI toggles, developer flags, transient prefs \\
\midrule
5. In-Memory Caches & \code{@Observable} properties & Session only &
   Derived state: PK curve, heat maps, adherence rate, screening log \\
\midrule
6. System State & UNNotificationCenter & System-managed &
   Pending reminder requests \\
\bottomrule
\end{tabular}
\caption{The six-tier persistence architecture. Tiers are ordered by
permissiveness: Tier~1 (Keychain) is the most protected; Tier~6
(UNNotificationCenter) is the least controlled by the application.}
\label{tab:tiers}
\end{table}

The architecture's central claim is that no data should live above its
sensitivity floor. Medication selection---the tuple of (hormone type, ester,
administration method)---directly identifies the user's medical regimen. It
resides in Tier~2 (encrypted UserDefaults). Reminder times---the hour and
minute at which a notification fires---carry no clinical information. They
reside in Tier~4 (plain UserDefaults). The tier assignment is explicit and
auditable: every UserDefaults key is listed in
Table~\ref{tab:defaults_keys} with its encryption status and rationale.

\subsection{Single Source of Truth per Manager}
\label{sec:single_source}

Each of the six managers owns its domain exclusively. No two managers write
to the same persistent key. The \code{mm.} prefix on
\code{MedicationManager}'s UserDefaults keys enforces this at the naming
level; the SwiftData model context is injected from \code{xTApp} rather than
fetched globally, ensuring the container's lifecycle is explicit and its
ownership unambiguous.

\begin{table}[h]
\centering
\small
\begin{tabular}{@{}p{2.8cm}p{3.2cm}p{3.0cm}p{5.5cm}@{}}
\toprule
\textbf{Manager} & \textbf{Files} & \textbf{Persistent Stores} &
\textbf{Key Responsibilities} \\
\midrule
\code{MedicationManager} & 7 files (main + 6 extensions) &
  SwiftData + encrypted UserDefaults &
  PK model, dose history, blood tests, calibration, site tracking \\
\midrule
\code{JournalManager} & 1 file & SwiftData &
  Mood entries, symptoms \\
\midrule
\code{ProfileManager} & 1 file & encrypted UserDefaults &
  Name, photo, date of birth \\
\midrule
\code{AuthenticationManager} & 1 file & plain UserDefaults &
  Biometric auth, auto-lock \\
\midrule
\code{NotificationManager} & 1 file & plain UserDefaults +
  UNNotificationCenter &
  Reminder scheduling \\
\midrule
\code{SecurityManager} & 1 file & Keychain &
  Encryption/decryption \\
\bottomrule
\end{tabular}
\caption{The six managers and their persistence responsibilities. All are
\code{@MainActor @Observable} except \code{SecurityManager}, which is
\code{Sendable}---a pure cryptographic service with no observable state.}
\label{tab:managers}
\end{table}

All six managers are injected via SwiftUI Environment in \code{xTApp.swift}
(lines~139--145). The injection is explicit: each manager is a
\code{static let shared} singleton, and the environment key's
\code{defaultValue} returns that shared instance. Views access managers
through \code{@Environment(\.medicationManager)} rather than through global
singleton access, which keeps the dependency graph visible and makes it
possible to inject test doubles in previews and unit tests.

\subsection{Persist at Mutation}
\label{sec:persist_at_mutation}

Mutation triggers persistence through a \code{didSet} pattern, not a
periodic save loop or a \code{scenePhase} hook. When
\code{selectedMedication} changes, \code{persistSelection()} fires
immediately (\code{MedicationManager+Persistence.swift}:74--80). When
\code{doseHistory} is modified, \code{saveContext()} commits to SwiftData
(\code{MedicationManager+Persistence.swift}:141--144). There is no window
between mutation and save where data can be lost on crash.

The pattern has a cost: every individual property mutation writes to
UserDefaults, which is a synchronous plist serialization. Rapid mutations
would amplify this cost. The \code{performBatch(\_:)} method
(\code{MedicationManager.swift}:178--186) addresses this by wrapping bulk
operations (e.g., seeding 110 sample doses) in a closure that suppresses
\code{refreshState()} until the batch completes, so the PK model is
recomputed once rather than 110 times:

\begin{verbatim}
func performBatch(_ work: () -> Void) {
    suppressRefresh = true
    work()
    suppressRefresh = false
    refreshState()
}
\end{verbatim}

The pattern is the persistence analog of the UX principle ``state facts,
don't advise'' (Section~2.1 of~\cite{xtux}): the manager states the mutation;
the \code{didSet} makes it durable. The view never asks for a save; the
manager never waits for one.

\subsection{Migrate, Never Lose}
\label{sec:migrate_never_lose}

Migration from unencrypted to encrypted storage is automatic and lossless.
The pattern---implemented identically in both \code{MedicationManager}
(\code{Persistence.swift}:146--149) and \code{ProfileManager}
(\code{ProfileManager.swift}:196--216)---is:

\begin{enumerate}[nosep,leftmargin=1.4em]
\item Try reading the encrypted key.
\item If absent, fall back to the unencrypted key.
\item If the unencrypted key succeeds, write the encrypted version and remove
  the plaintext.
\end{enumerate}

The migration is exercised on every cold launch via
\code{migrateToEncryptedStorage()}, which is idempotent: it checks for the
existence of the encrypted key before writing. A user who upgrades from a
build that stored medication selection in plain \code{selectedMedication} to a
build that stores it in \code{encryptedMedicationSelection} will have their
data transparently upgraded on first launch, with no data loss and no
prompt.

The migration also operates in the opposite direction for \emph{reading}:
\code{loadSelection()} (\code{Persistence.swift}:56--72) first attempts the
encrypted key; on miss, falls back to the legacy plaintext key. This means
the migration does not need to complete before the data is usable---the app
reads the legacy key if the migration has not yet run, and the migration
runs on the next write.

\subsection{Developer Safety Is Production Safety}
\label{sec:developer_safety}

Developer tooling must never destroy real user data. Two mechanisms enforce
this:

\textbf{Developer settings isolation.} All developer flags live under a
\code{dev.} prefix in UserDefaults (\code{DeveloperSettings.swift}:27--29).
In DEBUG builds, these flags are readable and writable. In RELEASE builds,
they are hardwired \code{false} and their setters are compiled out---no code
path in a RELEASE build can set a developer flag. On RELEASE launch,
\code{scrubStaleReleaseState()} (\code{DeveloperSettings.swift}:70--76)
removes every \code{dev.*} key left by a prior DEBUG install, preventing a
RELEASE-over-DEBUG install from inheriting stale developer state.

\textbf{Symbolic swap for sample profiles.} The developer sample profile
(``Jan'') seeds 110 synthetic dose records, blood tests, and journal entries
for testing. Enabling it must not destroy the user's real data.
\code{ProfileSnapshot} (DEBUG-only, \code{ProfileSnapshot.swift}) implements
a ``symbolic swap'': before Jan's data is seeded,
\code{snapshotIfNeeded()} captures every UserDefaults key and every
SwiftData record into a JSON file in the developer support directory
(\verb|xT-dev/user-snapshot.json|). When the sample profile is disabled,
\code{restore()} wipes the live store, reinserts the captured records, and
calls \code{rehydrateManagers()} to reset all in-memory state. The swap is
``symbolic'' because only records move---the CloudKit store itself is never
duplicated or replaced.

The \code{DosePayload} struct carries all 17 fields of \code{DoseRecord},
including the injection-site JSON and the dose-correction audit trail
(\code{isUserCorrected}, \code{originalDoseId},
\code{originalSiteJSON}). A comment at line~266 warns: ``any field
missing here is destroyed permanently the first time the user toggles the
Sample profile off.''

% ======================================================================
\section{The Persistence Architecture}
\label{sec:architecture}

This section walks through each layer of the persistence stack, tracing the
lifecycle of data from the moment the app launches to the moment it is
written to durable storage. It is the memory analog of the UX paper's ``User
Journey'' (Section~3 of~\cite{xtux}).

\subsection{Cold Launch Sequence}
\label{sec:cold_launch}

The cold-launch sequence in \code{xTApp.swift} establishes the persistence
stack in five steps:

\begin{enumerate}[nosep,leftmargin=1.4em]
\item \textbf{ModelContainer creation} (lines~54--87): A static
  \code{ModelContainer} is initialized with the schema of five
  \code{@Model} types and \code{cloudKitDatabase: .automatic}. If creation
  fails, the app falls back to an in-memory store---the on-disk store is
  never deleted, so a transient error (file corruption, lock contention) gets
  another chance on the next launch.
\item \textbf{Context wiring} (lines~91--92):
  \code{MedicationManager.shared.modelContext} and
  \code{JournalManager.shared.modelContext} receive the container's
  \code{mainContext}. The \code{modelContext} property's \code{didSet}
  triggers deferred loading: \code{JournalManager} calls
  \code{loadEntries()} when the context arrives.
\item \textbf{UserDefaults registration} (lines~97--106): All keys that are
  read but not yet written are registered with default values, preventing
  the ``found no value for key'' log noise that would otherwise accompany
  every cold launch before the user has configured reminders.
\item \textbf{Manager initialization}: Each manager's \code{init()} calls its
  \code{load*()} methods. UserDefaults-backed state loads synchronously.
  SwiftData-backed state loads when the context arrives via \code{didSet}.
\item \textbf{Derived state computation}: \code{MedicationManager.refreshState()}
  recomputes the PK curve, adherence rate, cycle position, and hormone level
  estimate from the loaded data.
\end{enumerate}

The sequence is deterministic and synchronous where possible. There is no
asynchronous loading spinner at launch; the data is either available
immediately (UserDefaults) or available within one run-loop turn
(SwiftData's mainContext fetch).

\subsection{SwiftData Schema}
\label{sec:swiftdata_schema}

\begin{table}[h]
\centering
\small
\begin{tabular}{@{}p{2.5cm}p{5.2cm}p{2.2cm}p{4.6cm}@{}}
\toprule
\textbf{Model} & \textbf{Key Fields} & \textbf{CloudKit} & \textbf{Notes} \\
\midrule
\code{DoseRecord} &
  17 stored properties: id, doseDate, medicationType, ester,
  administrationMethod, doseAmount, concentrationMgPerMl, notes, isDeleted,
  doseStatus, skipReason, injectionTime, injectionSitesJSON, recordedAt,
  isUserCorrected, originalDoseId, originalSiteJSON &
  Synced &
  Soft-delete via \code{isDeleted} flag; injection sites stored as raw JSON \\
\midrule
\code{JournalEntry} &
  id, date, mood, text, symptoms (JSON), hormoneLevel &
  Synced &
  Symptoms encoded as JSON array of strings \\
\midrule
\code{BloodTestRecord} &
  id, drawDate, measuredLevel, hormoneType, hoursSinceLastDose, notes,
  calibrationFactor (transient) &
  Synced &
  \code{calibrationFactor} is recomputed on load, not persisted \\
\midrule
\code{ScreeningRecord} &
  id, date, typeRawValue, notes &
  Synced &
  Screening type stored as raw String for forward compatibility \\
\midrule
\code{UserPreferences} &
  medicationRow, medicationColumn, lastDoseDate, nextDoseDate,
  concentrationDefaults, suppressionSpeed, spironolactoneClearanceEnabled,
  customFrequencyDays, defaultDoseAmount, defaultInjectionTime,
  showDoseIntensityInsights &
  Synced &
  One record per user; fetched as \code{FetchDescriptor.first} \\
\bottomrule
\end{tabular}
\caption{The five SwiftData \code{@Model} types. All participate in CloudKit
sync via \code{cloudKitDatabase: .automatic}.}
\label{tab:swiftdata_models}
\end{table}

The SwiftData schema was chosen over Core Data for three reasons: CloudKit
sync is a single configuration flag rather than a multi-step NSPersistentCloudKitContainer
setup; the \code{@Model} macro eliminates the need for separate
\code{.xcdatamodeld} files; and SwiftData's \code{FetchDescriptor} API is
native Swift rather than an Objective-C bridge.

\textbf{Soft-delete pattern.} \code{DoseRecord} carries an \code{isDeleted}
flag rather than being removed from the store. This preserves the audit trail
for dose corrections: a corrected dose retains its \code{originalDoseId}
reference to the original record, which must still exist in the store. The
\code{loadDoseHistory()} method (\code{Persistence.swift}:132--138) filters
out soft-deleted records:
\begin{verbatim}
doseHistory = (allRecords ?? []).filter { !$0.isDeleted }
\end{verbatim}

\textbf{CloudKit fallback.} If the on-disk \code{ModelContainer} fails to
initialize, the app falls back to \code{isStoredInMemoryOnly: true}
(\code{xTApp.swift}:78--86). The in-memory store survives the session but
not the launch. Critically, the code does \emph{not} delete the on-disk
store:
\begin{verbatim}
// Leave the on-disk store untouched so user data is never actively destroyed.
// A transient error gets another chance on the next launch.
\end{verbatim}
This means a corrupt store will produce session-only data loss rather than
permanent data destruction. The user is not explicitly notified---the app
launches normally, and data entered during the session is lost on next
launch---which is a known limitation (Section~\ref{sec:unaddressed}).

\subsection{UserDefaults with Encryption}
\label{sec:userdefaults}

\begin{table}[h]
\centering
\small
\begin{longtable}{@{}p{3.4cm}p{1.2cm}p{2.4cm}p{2.2cm}p{5.3cm}@{}}
\toprule
\textbf{Key} & \textbf{Enc?} & \textbf{Manager} & \textbf{Type} & \textbf{Purpose} \\
\midrule
\endhead
\code{encryptedMedicationSelection} & Yes & Medication & pipe-delimited String &
  Medication type $|$ method $|$ ester \\
\code{encryptedDoseDates} & Yes & Medication & JSON container &
  \code{\{lastDoseDate, nextDoseDate, hrtStartDate\}} \\
\code{encryptedProfileName} & Yes & Profile & UTF-8 String &
  User's chosen display name \\
\code{encryptedProfileImageData} & Yes & Profile & JPEG Data &
  Profile photo ($\leq$\,5\,MB) \\
\midrule
\code{profileDateOfBirth} & No & Profile & Date &
  Demographic, not health-identifiable \\
\code{mm.antiAndrogenInUse} & No & Medication & Bool &
  Whether an anti-androgen is configured \\
\code{mm.spironolactoneInUse} & No & Medication & Bool &
  Specific anti-androgen flag \\
\code{mm.spironolactoneClearanceEnabled} & No & Medication & Bool &
  Clearance adjustment toggle \\
\code{mm.suppressionSpeed} & No & Medication & String &
  Endogenous taper rate (fast/medium/slow) \\
\code{mm.customFrequencyDays} & No & Medication & Int? &
  User-set dosing interval override \\
\code{mm.missedDoseGracePeriodHours} & No & Medication & Double &
  Hours before a dose is marked missed \\
\code{mm.userTherapeuticTarget} & No & Medication & JSON data &
  \code{TherapeuticTarget} struct \\
\code{mm.avoidZonesJSON} & No & Medication & JSON data &
  \code{[String: AvoidZoneInfo]} dictionary \\
\midrule
\code{doseReminderEnabled} & No & Notification & Bool &
  Reminder toggle \\
\code{doseReminderHour} & No & Notification & Int &
  Hour of day (0--23) \\
\code{doseReminderMinute} & No & Notification & Int &
  Minute of hour (0--59) \\
\code{journalReminderEnabled} & No & Notification & Bool &
  Reminder toggle \\
\code{journalReminderHour} & No & Notification & Int &
  Hour of day (0--23) \\
\code{journalReminderMinute} & No & Notification & Int &
  Minute of hour (0--59) \\
\midrule
\code{requireAuthentication} & No & Auth & Bool &
  Face ID / passcode lock toggle \\
\code{autoLockTimeout} & No & Auth & TimeInterval &
  Grace period in seconds \\
\code{allowPasscodeFallback} & No & Auth & Bool &
  Whether to offer passcode after biometric failure \\
\midrule
\code{hasCompletedOnboarding} & No & ContentView & Bool &
  \code{@AppStorage} gate for main interface \\
\code{dev.*} (prefix) & No & Developer & Various &
  DEBUG-only; scrubbed on RELEASE \\
\bottomrule
\caption{Every UserDefaults key in the application, with encryption status
and rationale. The \code{mm.} prefix enforces \code{MedicationManager}'s
exclusive ownership of its transient preferences.}
\label{tab:defaults_keys}
\end{longtable}
\end{table}

The boundary between encrypted and unencrypted keys follows a simple rule:
\textbf{encrypt data that directly identifies the user's medical
regimen or personal identity; leave unencrypted data that carries
no clinical or identifying information.}

Medication selection is encrypted because the tuple \code{(``Estrogen'',
``Injectable'', ``Valerate (EV)'')} reveals the user's HRT regimen.
Reminder times are not encrypted because ``9:00\,AM'' carries no clinical
information. The \code{mm.*} transient preferences (anti-androgen flags,
suppression speed, therapeutic targets) are not encrypted because they are
numerical or boolean configuration values that, individually, do not
identify a medical regimen.

This is a judgment call. A user's spironolactone flag combined with their
custom frequency days and therapeutic targets \emph{could} reveal
clinical information in aggregate. The current design accepts this risk on
the reasoning that the combination requires physical device access, at which
point the attacker already has the device. A future iteration might move all
\code{mm.*} keys to encrypted storage; the migration path already exists.

\subsection{The Keychain and AES-GCM}
\label{sec:keychain}

\code{SecurityManager} (\code{SecurityManager.swift}, 97 lines) is a
\code{Sendable} singleton that manages the application's data-at-rest
encryption. It is the only manager that is not \code{@Observable}---it has
no observable state, only a cryptographic API.

\textbf{Key generation.} On first access, \code{getEncryptionKey()} generates
a 256-bit AES-GCM symmetric key via \code{SymmetricKey(size: .bits256)} from
Apple's CryptoKit. The key is serialized and stored in the Keychain under
service \code{com.starlacortes.xT.security}, account
\code{xTEncryptionKey}. If the key already exists in the Keychain, it is
fetched rather than regenerated.

\textbf{Keychain accessibility.} The key is stored with
\code{kSecAttrAccessibleWhenUnlockedThisDeviceOnly}. This attribute means:
\begin{itemize}[nosep,leftmargin=1.4em]
\item The key is only accessible when the device is unlocked (first
  unlock after boot is sufficient).
\item The key is \emph{not} included in unencrypted iTunes/Finder backups.
\item The key is \emph{not} synced to other devices via iCloud Keychain
  (\code{ThisDeviceOnly}).
\item The key may be evicted on device restore if the backup was not
  encrypted.
\end{itemize}

The choice of \code{ThisDeviceOnly} is deliberate: the encryption key should
never leave the device that generated it. But it creates the key-loss
problem exercised by scenario Nadia (Section~\ref{sec:scenarios}) and
documented as a structural limitation (Section~\ref{sec:unaddressed}).

\textbf{Encryption API.} \code{encryptData(\_:)} uses \code{AES.GCM.seal}
with a system-generated 96-bit nonce, returning the \code{combined}
representation (nonce $||$ ciphertext $||$ tag). \code{decryptData(\_:)}
reconstructs a \code{SealedBox} from the combined representation and calls
\code{AES.GCM.open}. The combined format is a standard: the nonce is
prepended, so each encryption is independent and there is no nonce-reuse
risk across encryptions.

\textbf{Error handling.} The \code{SecurityError} enum has five cases:
\code{keychainSaveFailed}, \code{keychainReadFailed},
\code{encryptionFailed}, \code{decryptionFailed}, and
\code{sealedBoxUnavailable}. The last is a safety net: the code uses
\code{guard let combined = sealedBox.combined} rather than a force-unwrap,
because ``health data is not worth a force unwrap'' (comment at
\code{SecurityManager.swift}:41).

\subsection{In-Memory Caches and Derived State}
\label{sec:in_memory}

The managers hold in-memory mirrors of persisted data for reactive SwiftUI
binding, plus derived state that is computed from that data. The caching
strategy distinguishes three categories:

\textbf{Direct mirrors} (loaded from store, mutated via manager methods):
\code{doseHistory}, \code{bloodTestHistory}, \code{journalEntries}. These
are loaded on init or context arrival and kept in sync with the store by
the \code{didSet} $\to$ \code{persist} pattern.

\textbf{Computed derived state} (recomputed by \code{refreshState()}):
\code{currentHormoneLevel}, \code{recentLevels}, \code{adherenceRate},
\code{onTimeStreak}, \code{regimenTransitionHistory}. These are pure functions
of the loaded data; they are recomputed whenever the underlying data changes.

\textbf{Cached computed state} with explicit invalidation:

\begin{itemize}[nosep,leftmargin=1.4em]
\item \code{heatMapCache} (\code{MedicationManager+Sites.swift}): A
  \code{@ObservationIgnored} dictionary mapping
  \code{"routeClass\_windowDays"} keys to per-zone injection-count
  dictionaries. Invalidated after each dose record via
  \code{invalidateHeatMapCache()}. Computation is offloaded to a
  \code{Task.detached(priority: .utility)} to avoid blocking the main actor.
\item \code{screeningLog} (\code{MedicationManager.swift}): A compact
  \code{[ScreeningType: Date]} lookup table, rebuilt from
  \code{screeningRecords}'s \code{didSet} observer. Marked
  \code{@ObservationIgnored} because it is only accessed programmatically,
  not rendered directly by SwiftUI views.
\end{itemize}

\textbf{Deliberately not memoized:} \code{calibrationFactor}. A comment at
\code{MedicationManager.swift}:259--262 explains: ``there is no mutation
hook a cache could hang invalidation off.'' The calibration factor depends on
blood test results, PK model constants, and dose history---any of which can
change independently. Rather than tracking every invalidation source, the
factor is recomputed on demand.

\textbf{Batch mutation guard.} The \code{suppressRefresh} flag and
\code{performBatch(\_:)} method prevent the quadratic cost of recomputing the
PK model after each of $n$ bulk dose insertions. During sample profile
seeding (110 doses), \code{suppressRefresh} is set to \code{true},
all doses are inserted, and \code{refreshState()} runs once at the end.

\subsection{File-System Resources}
\label{sec:file_system}

Two categories of data live on the file system rather than in SwiftData or
UserDefaults:

\textbf{Zone maps} (\code{ZoneMapFile.swift}). Baked per-vertex zone
identification data for the 3D body model is stored as JSON resources in the
app bundle (\verb|zone_map_<bodyType>.json|). Each zone map assigns one
\code{UInt16} zone ID to every vertex of the USDZ mesh ($\sim$14,000
vertices), encoded as a little-endian byte blob. The file carries a
\code{schemaVersion} field; \code{ZoneMap} refuses to load a file whose
version it does not recognize rather than silently mislabeling taps. Vertex
count is validated against the loaded mesh---a mismatch would scramble every
label. The previous design resolved taps by scanning 49 centroids; 57\% of
the mesh fell outside the centroid cutoffs. The per-vertex approach replaced
a heuristic with a lookup.

\textbf{Developer snapshot} (\code{ProfileSnapshot.swift}, DEBUG-only).
Captured to \verb|xT-dev/user-snapshot.json| in the application support
directory, deliberately outside the CloudKit-backed SwiftData store so
developer state can never sync. The snapshot contains all UserDefaults keys
(as a binary plist for lossless round-tripping) and all SwiftData records
(as Codable payloads). Restore wipes the live store, reinserts captured
records, and calls \code{rehydrateManagers()}.

\subsection{State Rehydration}
\label{sec:rehydration}

\code{ProfileSnapshot.rehydrateManagers()} (\code{ProfileSnapshot.swift}:196--247)
is the single code path that resets the application to a cold-launch state
without restarting the process. It:

\begin{enumerate}[nosep,leftmargin=1.4em]
\item Resets every \code{ProfileManager} property to empty, then calls
  \code{loadProfile()}.
\item Resets every \code{MedicationManager} property to its cold-launch
  default (nil medication, empty histories, nil hormone level), reloads
  transient preferences from UserDefaults, then calls
  \code{loadSelection()}, \code{loadDoseDates()},
  \code{loadDoseHistory()}, \code{loadBloodTestHistory()},
  \code{loadScreeningHistory()}, \code{refreshState()}, and
  \code{recomputeCalibrationFactors()}.
\item Resets \code{JournalManager}'s symptoms and calls
  \code{loadEntries()}.
\item Resets \code{NotificationManager}'s reminder settings from
  UserDefaults and calls \code{scheduleAllReminders()}.
\end{enumerate}

The rehydration sequence mirrors the cold-launch \code{init()} sequence
exactly, which means it exercises the same code paths. If
\code{rehydrateManagers()} works, cold launch works---there is no separate
``hot reload'' code path that could drift from the launch path.

\subsection{Developer State Isolation}
\label{sec:dev_isolation}

Developer settings live under a \code{dev.} prefix in UserDefaults
(\code{DeveloperSettings.swift}:27--29) and a dedicated support directory
(\verb|xT-dev/|). Three mechanisms protect production state from developer
tooling:

\begin{enumerate}[nosep,leftmargin=1.4em]
\item \textbf{Prefix isolation.} Every developer key carries the \code{dev.}
  prefix. \code{scrubStaleReleaseState()} iterates the UserDefaults
  dictionary and removes every key matching \code{dev.*}. No production key
  has this prefix.
\item \textbf{Compilation gating.} In RELEASE builds, the
  \code{developerModeEnabled} and \code{sampleProfileEnabled} properties are
  hardwired \code{let} constants set to \code{false}. Their \code{didSet}
  persistence code is compiled out. No code path in a RELEASE build can set
  a developer flag.
\item \textbf{Directory isolation.} The developer support directory
  (\verb|xT-dev/|) is outside the CloudKit-scoped SwiftData store, so
  snapshot files can never sync to iCloud---even if a DEBUG build were to
  accidentally leave a snapshot file in place.
\end{enumerate}

\subsection{Data Export}
\label{sec:data_export}

\code{DataExporter} (\code{DataExporter.swift}, 439 lines) is a stateless
formatter that produces portable CSV and JSON from live manager state. It
does not read raw stores; it aggregates state through an
\code{ExportState} struct that the caller populates from the managers.

The CSV export is multi-section: dose records, injection sites (one row per
site per dose), journal entries, blood tests, avoid zones, and user
preferences---each as a labeled section separated by blank lines. The JSON
export is a single \code{ExportPayload} with nested arrays for each record
type.

\textbf{Locale independence.} The export filename uses
\code{AppDateFormatters.isoDay}, which is POSIX- and UTC-locked. A device
on a Buddhist or Japanese calendar still writes
\verb|xT_doses_2026-08-04.csv| rather than
\verb|xT_doses_2569-08-04.csv|. The comment at line~249 explains: ``Filenames
must not vary by locale or calendar.''

\textbf{One-way export.} The \code{DataExporter} produces CSV and JSON but
has no corresponding \code{DataImporter}. A user who exports their data
cannot reimport it. This is a known limitation
(Section~\ref{sec:unaddressed}).

% ======================================================================
\section{Encoded Architecture Decisions}
\label{sec:encoded}

This section documents the persistence patterns, encryption boundaries, and
validation gates that ship in the application. It is the memory analog of the
UX paper's ``Encoded Design Decisions'' (Section~4 of~\cite{xtux}) and the PK
paper's ``Encoded Constants'' (Section~3 of~\cite{xtpk}): a reference for
what is in the codebase, where it came from, and what confidence attaches to
it.

\subsection{Persistence Patterns}
\label{sec:persistence_patterns}

\begin{table}[h]
\centering
\small
\begin{tabular}{@{}p{3.0cm}p{6.0cm}p{5.5cm}@{}}
\toprule
\textbf{Pattern} & \textbf{How It Works} & \textbf{Where} \\
\midrule
\code{didSet} $\to$ persist &
  Property observer calls a \code{persist*()} method on mutation &
  \code{MedicationManager.swift},
  \code{AuthenticationManager.swift} \\
\code{load} $\to$ hydrate &
  \code{init()} calls \code{load*()} methods; SwiftData load deferred to
  \code{modelContext} \code{didSet} &
  All managers \\
\code{encrypt} $\to$ envelope &
  \code{SecurityManager.encryptData()} wraps value in AES-GCM combined blob &
  \code{SecurityManager.swift} \\
Auto-migration &
  On load, try encrypted key; on miss, try legacy plaintext; on hit,
  write encrypted, delete plaintext &
  \code{MedicationManager+Persistence.swift}:146,
  \code{ProfileManager.swift}:196 \\
SwiftData \code{FetchDescriptor} &
  \code{FetchDescriptor<Model>(sortBy: [...])} $\to$
  \code{context.fetch()} &
  \code{Persistence.swift}:132--138 \\
Soft delete &
  Set \code{isDeleted = true}; filter on load &
  \code{DoseRecord.swift}:31,
  \code{Persistence.swift}:137 \\
Deferred load &
  \code{modelContext} \code{didSet} triggers
  \code{loadEntries()} or equivalent &
  \code{JournalManager.swift}:21,
  \code{MedicationManager.swift}:117--122 \\
Symbolic swap &
  Snapshot all records to JSON; wipe live store; reinsert from JSON;
  rehydrate managers &
  \code{ProfileSnapshot.swift}:141--191 \\
Batch guard &
  \code{suppressRefresh = true}; perform work;
  \code{suppressRefresh = false}; \code{refreshState()}
  once &
  \code{MedicationManager.swift}:178--186 \\
\bottomrule
\end{tabular}
\caption{The nine persistence patterns used across the codebase. Each
pattern has a single canonical implementation; no pattern is reinvented in a
second location.}
\label{tab:patterns}
\end{table}

\subsection{Encryption Boundaries}
\label{sec:encryption_boundaries}

\begin{table}[h]
\centering
\small
\begin{tabular}{@{}p{3.8cm}p{1.8cm}p{8.9cm}@{}}
\toprule
\textbf{Data} & \textbf{Encrypted?} & \textbf{Rationale} \\
\midrule
Medication type, ester, method & Yes &
  Directly identifies HRT regimen \\
Dose dates (last, next, HRT start) & Yes &
  Reveals treatment timeline and start date \\
Profile name & Yes &
  May identify user personally \\
Profile photo & Yes &
  Face image, highest sensitivity \\
Date of birth & No &
  Demographic, not uniquely identifying \\
Reminder times & No &
  Clock times carry no clinical information \\
Anti-androgen flag & No &
  Binary toggle, low individual sensitivity \\
Suppression speed & No &
  ``fast''/``medium''/``slow'' enum value \\
Custom frequency days & No &
  Integer; individual sensitivity low \\
Therapeutic targets & No &
  User-set numerical ranges \\
Dose records & No (CloudKit) &
  Structural data; CloudKit encrypts in transit and at rest \\
Journal entries & No (CloudKit) &
  User-authored content \\
Blood test results & No (CloudKit) &
  Lab values; CloudKit encrypts in transit and at rest \\
\bottomrule
\end{tabular}
\caption{Encryption status for every category of persisted data. ``No
(CloudKit)'' means the data relies on CloudKit's built-in encryption rather
than application-layer AES-GCM.}
\label{tab:encryption}
\end{table}

The encryption boundary is a conscious split. SwiftData records (dose
history, journal entries, blood tests) are not additionally encrypted by the
application because: (a) CloudKit encrypts data in transit (TLS) and at rest
(on Apple's servers); (b) the records are structural and benefit from
CloudKit's sync, query, and merge capabilities; (c) adding an
application-layer encryption envelope over SwiftData records would prevent
CloudKit from performing server-side merge operations. The trade-off is that
Apple holds the CloudKit encryption keys, not the user---a limitation
discussed in Section~\ref{sec:unaddressed}.

\subsection{Validation Gates}
\label{sec:validation_gates}

\begin{table}[h]
\centering
\small
\begin{tabular}{@{}p{3.2cm}p{2.4cm}p{4.2cm}p{4.7cm}@{}}
\toprule
\textbf{Gate} & \textbf{Value} & \textbf{File} & \textbf{Line} \\
\midrule
Profile image max size & 5\,MB &
  \code{ProfileManager.swift} & 30 \\
Profile name max length & 50 chars &
  \code{ProfileManager.swift} & 27 \\
JPEG compression quality & 0.8 &
  \code{ProfileManager.swift} & 24 \\
Calibration factor $\alpha$ clamp & $[0.2, 5.0]$ &
  \code{MedicationManager+Dosing.swift} & --- \\
Schema version check & 1 &
  \code{ZoneMapFile.swift} & 28 \\
Vertex count validation & matches loaded mesh &
  \code{ZoneMapFile.swift} & 66--68 \\
Dose volume warning & $> 3.0$\,mL &
  \code{MedicationManager+Dosing.swift} & --- \\
Injection dose sanity & $\leq 1000$\,mg &
  \code{MedicationManager+Dosing.swift} & --- \\
\bottomrule
\end{tabular}
\caption{Input validation gates that prevent silent data corruption. Values
outside these ranges are rejected with an explicit error or warning rather
than silently persisted.}
\label{tab:gates}
\end{table}

\subsection{Data Flow}
\label{sec:data_flow}

The reactive data pipeline flows in one direction, with no cycles:

\begin{verbatim}
User Action
    |
    v
@Observable property mutation (Manager)
    |
    +---> didSet ---> persist*() ---> SwiftData / UserDefaults
    |
    +---> refreshState() (if not suppressed)
              |
              v
         Derived State (currentHormoneLevel, recentLevels,
                        adherenceRate, regimenHistory)
              |
              v
         SwiftUI Views (re-render via @Observable)
\end{verbatim}

Views observe manager state through \code{@Environment} and SwiftUI's
\code{@Observable} tracking. They never call \code{save()} or
\code{load()}---the manager is the API, and persistence is an implementation
detail of property mutation. A view that sets
\code{medicationManager.selectedMedication = "Estrogen"} does not know or
care that this triggers AES-GCM encryption and a UserDefaults write; it
mutates a property and the manager handles the rest.

\subsection{Keychain Accessibility}
\label{sec:keychain_design}

\begin{table}[h]
\centering
\small
\begin{tabular}{@{}p{3.2cm}p{5.5cm}p{5.8cm}@{}}
\toprule
\textbf{Attribute} & \textbf{Value} & \textbf{Effect} \\
\midrule
Class & \code{kSecClassGenericPassword} &
  Standard Keychain item type for arbitrary data \\
Service & \code{com.starlacortes.xT.security} &
  Namespaced to the app's bundle identifier \\
Account & \code{xTEncryptionKey} &
  Distinguishes the encryption key from other Keychain items \\
Accessibility & \code{kSecAttrAccessibleWhenUnlockedThisDeviceOnly} &
  Not accessible when device is locked; not included in unencrypted
  backups; not synced via iCloud Keychain \\
\bottomrule
\end{tabular}
\caption{Keychain attributes for the AES-256-GCM encryption key. The
\code{ThisDeviceOnly} suffix is load-bearing: it prevents the key from
appearing on other devices signed into the same iCloud account.}
\label{tab:keychain}
\end{table}

% ======================================================================
\section{Validation: Persistence Scenarios}
\label{sec:scenarios}

Section-level descriptions of persistence layers describe the architecture
as designed. Scenarios describe the architecture as \emph{exercised}---by
specific failure modes, edge cases, and transitions. This section adopts the
scenario-validation approach of the companion UX paper (Section~5
of~\cite{xtux}), exercising memory-architecture edge cases rather than UX
edge cases.

\begin{longtable}{@{}p{2.6cm}p{5.0cm}p{6.9cm}@{}}
\toprule
\textbf{Scenario} & \textbf{Event} & \textbf{What it tests, and the result} \\
\midrule
\endhead
Maya, cold launch &
  Fresh install, no prior data. Completes onboarding; force-quits; relaunches.
  & \textbf{Cold-launch persistence.} \code{MedicationManager.init()} loads
    from empty UserDefaults. Onboarding writes medication selection, dose
    dates, and user preferences. On restart, all values are decrypted and
    loaded. SwiftData container initializes with CloudKit schema. The PK
    curve shows the single post-onboarding dose. \good{Pass.} \\[0.4em]

Dev, legacy migration &
  Upgrading from a build that stored medication selection unencrypted in the
  \code{selectedMedication} key.
  & \textbf{Unencrypted-to-encrypted migration.}
    \code{migrateToEncryptedStorage()} checks for
    \code{encryptedMedicationSelection}; on miss, \code{loadSelection()}
    falls back to \code{selectedMedication}; on hit, writes encrypted version
    and removes plaintext. Both \code{MedicationManager} and
    \code{ProfileManager} follow identical paths. \good{Pass.} \\[0.4em]

Priya, iCloud sync conflict &
  Enters a dose on iPhone; later opens iPad. Both devices have pending
  CloudKit changes for the same record.
  & \textbf{CloudKit merge.} SwiftData's \code{.automatic} CloudKit database
    handles merge. \code{DoseRecord}'s \code{@Attribute(.unique) var id}
    prevents duplication. The \code{isDeleted} soft-delete flag survives
    merge conflicts---a record deleted on one device stays deleted after
    sync. \good{Pass (delegated to SwiftData).} \\[0.4em]

Sam, ModelContainer failure &
  File corruption or schema version mismatch prevents SwiftData container
  creation.
  & \textbf{Graceful degradation.} \code{xTApp.swift}:78--86 creates an
    in-memory store. The app launches normally. Dose records entered during
    the session are lost on next launch, but UserDefaults data (medication
    selection, profile) survives. On next launch with cleared corruption,
    SwiftData reinitializes from CloudKit or a fresh store. \good{Pass
    (graceful degradation).} \\[0.4em]

Rowan, sample profile swap &
  Developer toggles Sample profile on (Jan's data is seeded), uses it,
  toggles off.
  & \textbf{Symbolic swap integrity.}
    \code{ProfileSnapshot.snapshotIfNeeded()} captures all 68+ UserDefaults
    keys and all SwiftData records before Jan's data is seeded. On
    toggle-off, \code{restore()} wipes all live records, reinserts captured
    ones, and calls \code{rehydrateManagers()}. The user's original data is
    byte-identical to pre-toggle state. The \code{DosePayload} carries all
    17 fields. \good{Pass.} \\[0.4em]

Kai, RELEASE over DEBUG &
  User installs App Store build over a TestFlight DEBUG build that had
  developer flags set (\code{dev.developerModeEnabled = true}).
  & \textbf{Developer state scrub.} \code{xTApp.init()} calls
    \code{scrubStaleReleaseState()} in RELEASE. All \code{dev.*} keys are
    removed. The \verb|xT-dev/| directory is deleted.
    \code{developerModeEnabled} is hardwired \code{false}. No stale state
    survives. \good{Pass.} \\[0.4em]

Jules, data export &
  User exports all data via Settings $\to$ Data Export after 6 months of
  use.
  & \textbf{Export completeness.} \code{DataExporter} reads live manager
    state. CSV contains 6 labeled sections: dose records (with injection
    sites as separate rows), journal entries, blood tests, avoid zones, and
    preferences. JSON includes all fields from all model types including
    the dose-correction audit trail. Filename uses POSIX-locked date
    formatter. Export is one-way---there is no import path. \good{Pass
    (export only).} \\[0.4em]

Nadia, encryption key loss &
  Device is restored from an unencrypted backup. The Keychain item is
  evicted. The app launches.
  & \textbf{Key eviction behavior.}
    \code{SecurityManager.getEncryptionKey()} fails to read from Keychain,
    generates a new key. The new key cannot decrypt existing
    \code{encryptedMedicationSelection}, \code{encryptedDoseDates},
    \code{encryptedProfileName}, or \code{encryptedProfileImageData}.
    \code{loadSelection()} falls back to legacy unencrypted keys (which
    don't exist post-migration), loads empty defaults. The app launches
    with no medication configured---the user must re-enter their
    medication details and profile. Dose records and journal entries
    (SwiftData/CloudKit) survive. \bad{Fail with data loss.} \\
\bottomrule
\end{longtable}

Scenario Nadia is the only failure in the set. It exercises the
symmetric-key model's structural limitation: a single key protects all
encrypted UserDefaults values, and losing that key means losing all of them.
The scenario is included because the UX paper's principle of ``honest
presentation'' (Section~2.1 of~\cite{xtux}) applies equally to architecture
documentation: limitations must be named, not hidden.

% ======================================================================
\section{Discussion: Design Trade-offs and Limitations}
\label{sec:discussion}

This section describes trade-offs that were consciously accepted and
limitations that remain unaddressed. It is the memory analog of the UX
paper's Discussion (Section~6 of~\cite{xtux}) and the PK paper's
Discussion (Section~5 of~\cite{xtpk}).

\subsection{The Split-Persistence Trade-off}
\label{sec:split_persistence}

Medication selection and dose dates live in encrypted UserDefaults
(Tier~2), while dose records and blood tests live in SwiftData (Tier~3).
This split means that a complete picture of the user's medication
regimen requires reading from two stores with different encryption
properties, different query capabilities, and different sync behavior.

The split is justified by three factors:

\begin{enumerate}[nosep,leftmargin=1.4em]
\item \textbf{Access frequency.} Medication selection is read once per
  launch to determine which PK constants to use. It changes only when the
  user changes their medication in Settings---a rare event. UserDefaults is
  optimized for this pattern: read-once, write-rarely, key-value.
\item \textbf{Sensitivity.} The tuple (hormone type, ester, method) directly
  identifies the user's medical regimen. Individual dose records are less
  revealing: a single record of ``4\,mg EV, 2026-03-15'' reveals a dose
  amount but not the regimen structure. The sensitivity difference justifies
  the encryption difference.
\item \textbf{Sync requirements.} Dose records benefit from cross-device
  sync---a user should see their full dose history on both iPhone and iPad.
  Medication selection is unlikely to change between devices and does not
  benefit from sync. Encrypting data that needs CloudKit merge would be
  counterproductive: CloudKit cannot merge encrypted blobs.
\end{enumerate}

The cost is architectural complexity: a developer reading the codebase must
understand that ``medication data'' is split across two stores, and a change
to the medication model may require changes to both.

\subsection{DidSet Persistence vs.\ Periodic Save}
\label{sec:didSet_tradeoff}

The \code{didSet} $\to$ \code{persist()} pattern writes to storage on every
mutation. This trades write frequency for safety: no data is ever lost
between the mutation and a periodic save. The cost is that rapid mutations
would trigger excessive writes, which is why \code{performBatch(\_:)}
suppresses \code{refreshState()} during bulk operations.

An alternative design would use a periodic save (e.g., on
\code{scenePhase} \code{.background}) with an in-memory change queue. This
would reduce write frequency at the cost of a data-loss window between the
last mutation and the next background event. For a health-tracking
application where a dose record is the most important user action, that
window is unacceptable: a user who records a dose and immediately
force-quits must find the dose present on relaunch.

\subsection{Symmetric Key in Keychain vs.\ Data Protection API}
\label{sec:symmetric_vs_dp}

The current design generates a single 256-bit AES-GCM key and stores it in
the Keychain. An alternative would use Apple's Data Protection API
(\code{NSFileProtectionComplete}), which encrypts files at the OS level
without application-managed keys. The symmetric-key approach was chosen
because: (a) it works independently of file-level protection, allowing the
app to encrypt individual UserDefaults values rather than entire files;
(b) it allows selective encryption---some UserDefaults keys are encrypted,
others are not; (c) the Keychain accessibility attribute
\code{WhenUnlockedThisDeviceOnly} provides granular control over when the
key is available.

The cost is key management complexity and the key-loss problem (scenario
Nadia, Section~\ref{sec:scenarios}).

\subsection{SwiftData + CloudKit vs.\ Local-Only Store}
\label{sec:swiftdata_vs_local}

SwiftData with CloudKit was chosen over a local-only store (Core Data,
raw SQLite) because it provides cross-device sync without a server. The
trade-offs: (a) CloudKit sync introduces latency and potential merge
conflicts; (b) the schema is constrained by CloudKit's type system (enums
stored as raw String values throughout the codebase, e.g.,
\code{ScreeningRecord.typeRawValue}); (c) the app is tied to the user's
iCloud account---a user without iCloud gets no sync.

The benefit is that no server infrastructure exists, no user accounts are
created, and no data leaves the user's iCloud account. For a health
application, this is the right trade: ``no server'' means no server to
breach.

\subsection{What Remains Unaddressed}
\label{sec:unaddressed}

The following limitations are structural---they cannot be addressed by
refining the existing architecture and would require new features or
infrastructure.

\textbf{No end-to-end encryption of SwiftData.} CloudKit encrypts data in
transit (TLS) and at rest on Apple's servers, but Apple holds the encryption
keys. The application does not add an encryption layer over SwiftData
records. A user who wants full end-to-end encryption of all health data---so
that not even Apple can access it---would need a different architecture
(application-layer encryption over each record, with key material that never
touches Apple's servers).

\textbf{Key recovery on device restore.} If a user restores their device from
a backup that does not include Keychain items (iTunes backup without
encrypted backup enabled), the encryption key is lost. Encrypted UserDefaults
values become permanently unreadable. The app does not surface this risk or
offer a recovery mechanism. Mitigations exist: encrypted iTunes backups
preserve Keychain items; iCloud backups always preserve them. But the app
does not guide the user toward these mitigations.

\textbf{No data integrity verification.} The app does not compute checksums
or MACs over persisted data. A corrupted UserDefaults plist or SwiftData
store would produce silent errors (failed decryption, zero-length arrays)
rather than explicit corruption warnings. Adding integrity verification
would require a MAC over each persisted value, which would need its own key
management---the current single-key model would need to become a two-key
model (one for confidentiality, one for integrity).

\textbf{No backup/restore beyond iCloud.} The \code{DataExporter} produces
CSV and JSON but has no corresponding \code{DataImporter}. A user who
exports their data cannot reimport it. The export is a one-way safety hatch
for the user's own records, not a migration tool.

\textbf{No selective sync.} All SwiftData models participate in CloudKit
sync. There is no mechanism for a user to say ``sync dose records but keep
journal entries local-only.'' The all-or-nothing sync model is simpler but
less flexible.

\textbf{Silent fallback on decryption failure.} When
\code{loadSelection()} fails to decrypt the encrypted medication selection,
it falls back to the legacy unencrypted key; if that is also absent, it
loads empty defaults. The user is not notified that their encrypted data
could not be read. This is a deliberate choice: surfacing a ``decryption
failed'' error on launch would be alarming to a user who does not know what
encryption is, and the most likely cause (key eviction on restore) is not
something the user can fix from within the app. But it means data loss can
occur silently.

\textbf{Developer safety is DEBUG-only.} The \code{ProfileSnapshot}
mechanism that protects user data during sample profile testing is compiled
out of RELEASE builds. If a RELEASE build ever needed a similar safety net
(e.g., for a data repair tool or a migration assistant), the infrastructure
would not exist.

\subsection{Architecture Confidence Hierarchy}
\label{sec:confidence}

Analogous to the UX paper's design decision confidence hierarchy
(Section~6.7 of~\cite{xtux}) and the PK paper's source quality hierarchy
(Section~5.3 of~\cite{xtpk}), persistence decisions in xT rest on a
hierarchy of confidence:

\begin{enumerate}[nosep,leftmargin=1.4em]
\item \textbf{Platform-standard} (highest confidence). Decisions that use
  Apple frameworks as designed: SwiftData \code{@Model} persistence,
  CloudKit \code{.automatic} sync, Keychain Services API,
  CryptoKit AES-256-GCM, \code{UNUserNotificationCenter},
  \code{LocalAuthentication}. These are tested across billions of devices.
\item \textbf{Security-grounded.} Decisions that follow from cryptographic
  best practices: on-device key generation, AES-GCM with system-generated
  nonces, \code{kSecAttrAccessibleWhenUnlockedThisDeviceOnly}. These are
  the standard recommendations for data-at-rest encryption on iOS.
\item \textbf{Reasoned but untested.} Decisions based on architectural
  reasoning without formal performance or security testing: the
  \code{didSet} persistence pattern's behavior under rapid mutations, the
  encrypted UserDefaults key count's impact on cold-launch time, the
  \code{ProfileSnapshot}'s memory footprint with large dose histories
  (1,000+ records), the CloudKit merge behavior for
  \code{InjectionSite} JSON blobs. These are internally consistent and
  follow from the design principles, but their edge-case behavior has not
  been profiled.
\item \textbf{Placeholder} (lowest confidence). Decisions known to be
  suboptimal but retained because a better alternative is not yet available:
  the lack of key recovery on device restore, the absence of integrity
  verification, the one-way data export, the silent fallback on decryption
  failure. These are documented gaps; closing them would require
  architectural changes rather than incremental improvements.
\end{enumerate}

% ======================================================================
\section{Conclusion}
\label{sec:conclusion}

We have presented the memory architecture of xT, a mobile application for
hormone replacement therapy monitoring. The architecture's key structural
features are:

\begin{itemize}[nosep,leftmargin=1.4em]
\item A \textbf{six-layer persistence stack} that assigns every piece of data
  to the least-permissive tier that can serve its access pattern: Keychain
  for the encryption key, encrypted UserDefaults for health-identifiable
  configuration, SwiftData with CloudKit for structured longitudinal data,
  plain UserDefaults for non-sensitive preferences, in-memory caches for
  derived state, and UNNotificationCenter for system-level reminders.
\item An \textbf{encryption model} in which a single AES-256-GCM symmetric
  key---generated on-device, stored in the Keychain with
  \code{WhenUnlockedThisDeviceOnly}, and never transmitted---protects the
  four health-identifiable UserDefaults values, with automatic migration of
  legacy unencrypted data.
\item A \textbf{manager architecture} of six \code{@MainActor @Observable}
  singletons, each owning its domain exclusively and persisting via a
  \code{didSet}-driven pattern that eliminates the save-to-crash window, with
  a \code{performBatch} guard for bulk operations and documented
  invalidation strategies for cached derived state.
\item A \textbf{developer safety system} in which \code{dev.*}-prefixed
  UserDefaults keys and a dedicated support directory isolate developer
  state from production data, \code{ProfileSnapshot} implements a symbolic
  swap to protect user data during sample profile testing, and
  \code{scrubStaleReleaseState()} guarantees that a RELEASE build never
  inherits developer flags.
\item A \textbf{rehydration mechanism} (\code{ProfileSnapshot.rehydrateManagers()})
  that resets the application to cold-launch state without restarting the
  process, exercising the same load paths as a genuine cold launch.
\end{itemize}

The eight persistence scenarios (Section~\ref{sec:scenarios}) exercise the
architecture across cold launch, migration, sync conflicts, storage failure,
developer safety, data export, and key loss. Seven pass; one (key loss on
device restore) fails with data loss, which is named as a structural
limitation rather than hidden.

Design trade-offs are stated plainly: the split between encrypted
UserDefaults and SwiftData reflects a sensitivity/access-frequency
judgment; the \code{didSet} persistence pattern prioritizes safety over write
efficiency; the symmetric-key model provides selective encryption at the cost
of key management complexity; the CloudKit dependency provides server-free
sync at the cost of reduced control over encryption key custody.

What remains unaddressed---end-to-end encryption of SwiftData records, key
recovery on device restore, data integrity verification, backup/restore
beyond iCloud, selective sync, and silent fallback on decryption
failure---is listed without apology.

The most important property of a health-data persistence architecture, in
our view, is the same as the most important property of the pharmacokinetic
model that drives it and the user interface that delivers it: \textbf{honesty
about what it can and cannot protect.} xT's memory architecture is designed
to be honest. It encrypts what must be encrypted and documents what it does
not. It names its key-loss scenario rather than pretending it does not exist.
It states which trade-offs are conscious and which are inherited. This
honesty---the explicit marking of what is protected, what is not, and
why---is not the most visible property of the persistence layer, but it is
the one that matters most for an application whose users trust it with their
medical data.

\vfill
\noindent\rule{\linewidth}{0.4pt}\\[0.3em]
{\footnotesize This paper describes a mobile application's persistence
architecture. It is not a security audit, and the architecture decisions
described here should not be interpreted as security guarantees. The
application's encryption model protects data at rest on the device but does
not provide end-to-end encryption of CloudKit-synced records. The
symmetric-key model has no recovery mechanism if the Keychain item is
evicted.}

% ======================================================================
\begin{thebibliography}{99}

\bibitem{xtux}
Cortes S.\ (2026). ``Progressive Clarity with Honest Constraint: A Design
Framework for Hormone Replacement Therapy Tracking in the xT Mobile
Application.'' August 2026.

\bibitem{xtpk}
Cortes S.\ (2026). ``Superposed Bateman Biexponentials with
Clearance-Constrained Scaling and Per-User Calibration: A Pharmacokinetic
Framework for Hormone Replacement Therapy Monitoring in the xT Mobile
Application.'' August 2026.

\bibitem{apple_cryptokit}
Apple Inc.\ (2024). ``CryptoKit.'' Apple Developer Documentation.
\verb|https://developer.apple.com/documentation/cryptokit|

\bibitem{apple_security}
Apple Inc.\ (2024). ``Keychain Services.'' Apple Developer Documentation.
\verb|https://developer.apple.com/documentation/security/keychain_services|

\bibitem{apple_swiftdata}
Apple Inc.\ (2024). ``SwiftData.'' Apple Developer Documentation.
\verb|https://developer.apple.com/documentation/swiftdata|

\bibitem{apple_cloudkit}
Apple Inc.\ (2024). ``CloudKit.'' Apple Developer Documentation.
\verb|https://developer.apple.com/documentation/cloudkit|

\bibitem{apple_security_guide}
Apple Inc.\ (2024). ``Apple Platform Security.'' Apple Platform Security
Guide. \verb|https://help.apple.com/pdf/security/en_US/apple-platform-security-guide.pdf|

\end{thebibliography}

\end{document}
